\section{Generate bytes before transforming a message}
A stream cipher produces a sequence of bytes called a \emph{keystream}. Encryption
combines each message byte with the corresponding keystream byte using exclusive OR.
Decryption performs the same operation. If a message byte is \texttt{4c} and the
keystream byte is \texttt{22}, the encrypted byte is \texttt{6e}. Applying XOR with
\texttt{22} again recovers \texttt{4c}. The byte's position matters: decrypting with
the next keystream byte generally gives a different result.

This project builds that process from a four-word calculation. First trace a
quarter-round and check its intermediate values. Then apply quarter-rounds to selected
positions in a sixteen-word state, producing one 64-byte block. After the scalar block
works, place four independent quarter-rounds in the four word lanes of MSA registers.
Finally, use complete blocks and a partial last block to transform a message.

The first observation is deliberately smaller than an encrypted message. If four
words are wrong after a quarter-round, inspecting the file format or the XOR loop
cannot repair them. Likewise, if a complete block agrees but a 65-byte message fails,
inspect the transition to the next block and the final one-byte store. The boundaries
in this chapter give each failure a short path back to the calculation that produced it.

The algorithm here is the IETF form of ChaCha20, with an eight-word key, a one-word
counter and a three-word nonce. A nonce identifies an encryption under a key; the
counter identifies a block within that encryption. This exercise uses explicit test
values so results can be reproduced. It does not provide key generation, nonce
allocation, a message authentication tag or a deployable encryption format.
